\documentclass[11pt]{article}
\usepackage[T1]{fontenc}
\usepackage[utf8]{inputenc}
\usepackage{amsmath,amssymb,amsthm}
\usepackage[margin=1in]{geometry}
\usepackage{hyperref}
\usepackage{lmodern}
\setlength{\emergencystretch}{3em}
\title{P02.3: Weakly Smooth Decreasing Regularization, Singular Curves, and an Analytic Boundary Obstruction}
\author{Unverified research draft}
\date{4 October 2026}
\begin{document}
\maketitle

This is an unverified research draft. The full characterization for arbitrary classical psh germs on higher-dimensional good strict analytic spaces remains open. The results below are sufficient sectors and a sharp obstruction within a particular boundary family. Successful source checks and TeX builds are not mathematical verification.

\section{Conventions and checked published premises}

An analytic extension means a complete isometric extension of the ground field. All neighborhoods and overlaps in our conclusions are ordinary open sets. Classical psh means GR, Definition 4.1: an upper semicontinuous function with values in \(\mathbb R\cup\{-\infty\}\) whose pullback to every analytic curve after every analytic extension is subharmonic. Identically minus infinity is permitted. For smooth connected strict curves over a nontrivially valued field, GR 3.4.1 identifies this class with Thuillier's subharmonic functions together with the constant minus infinity. No continuity or finite-value hypothesis is imposed on the starting function.

We use the following checked primary results. References refer to the cached versions, not to an unpublished book.

GR denotes Gubler--Rabinoff, arXiv:2506.13548v2. Proposition 3.5(3) says that an upper semicontinuous function on a curve which is subharmonic on its analytic interior is subharmonic on the whole curve. Proposition 4.2 gives ordinary sheafiness, positive linear combinations and maxima, decreasing nets, scalar extension in both directions, pullback, and descent along finite surjective morphisms. Theorem 6.1 gives the local maximum principle at analytic interior points. Theorem 6.5 gives stability under a pointwise convergent net when its limit is upper semicontinuous. Lemma 3.9 describes finite normalization of reduced strict curves: over a perfect field it is rig-smooth, is an isomorphism outside discrete rigid points, is a quotient, and locally splits into finitely many components, each meeting exactly one normalization preimage of the singular point. Its proof also records local constancy of invertible analytic absolute values at rigid points.

GJR denotes Gubler--Jell--Rabinoff, arXiv:2111.05741. Its standing field is complete and nontrivially valued and its spaces are good and strictly analytic. Definition 5.1 uses the G-topology generated by strict affinoid domains. Definition 8.1 defines harmonic tropicalizations using integer-PL harmonic coordinates. Proposition 7.5 expresses real harmonic germs as real linear combinations of integer harmonic germs. Proposition 7.6 and Remark 8.2(4) give pullback and arbitrary scalar extension of harmonic coordinates. Definitions 10.3 and 10.7 define weakly smooth functions by smooth Lagerberg functions pulled back along harmonic tropicalizations, with sheafification in the ordinary topology; 10.8 gives pullback. Proposition 10.9 describes common-chart equality. Remark 10.10 includes ordinary CLD forms in weakly smooth forms. Proposition 14.2 gives finite Galois descent for weakly smooth forms. Proposition 14.5(1) states equality of integrals under finite Galois scalar extension, without a degree factor. Sections 11 and 12 provide integration and strong currents on Hausdorff boundaryless pure-dimensional spaces, with Stokes and \(d'd''[v]=[d'd''v]\) for weakly smooth functions.

GJR Proposition 15.7 identifies real harmonic functions on rig-smooth curves with Thuillier harmonic functions. Proposition 15.9 concerns piecewise \((\mathbb Z,\Gamma)\)-linear harmonic skeleton functions on the indicated strict subgraphs. It is not a statement about arbitrary real bends on arbitrary domains. Remark 15.10 allows harmonic functions which are weakly smooth but not ordinary CLD smooth. \textbf{Proposition 15.11 states a decreasing net}, on each relatively compact ordinary open subset of a smooth boundaryless curve, for arbitrary Thuillier-subharmonic targets. The proof's last sentence extracts a sequence for an individual piecewise real-linear target; it does not replace the preceding arbitrary-target net by a sequence.

Thuillier, thesis tel-00010990v1, distinguishes the all-domain site \(X_G\) from the strict rigid site \(X_R\) in 2.1.1. Definition 3.2.1 of lisse uses all-domain affinoid coverings. Theorem 3.4.2 and its proof give a decreasing family indexed by finite subsets, not a sequence. We do not assert strict GJR real-PL membership for every such lisse function.

Wanner, arXiv:1801.04713v3, assumes an algebraically closed complete nontrivially valued field. Corollary 5.4 assumes a smooth proper algebraic curve, residue field algebraic over a finite field or a projective-line/Mumford-curve alternative, and a finite continuous current-psh target on an open subset. Its compact Dini argument concerns that continuous target. It is not an arbitrary classical-potential theorem. Only the elementary regularized-maximum mechanism in its Lemma 5.2 and the proof of GJR 15.11 is used here.

CD denotes the actual saved arXiv:1204.6277v2. Its relevant numbering is Definition 8.1.1, Lemma 8.1.4, Proposition 8.1.5, and Proposition 8.3.1, for ordinary smooth psh generators, transfinite stages, pullbacks and current positivity. Its analytic ordinary-smooth description is Corollary 15.1.2; Lemma 1.2.11 gives linear pullback of positivity. These statements do not identify all classical functions with the first smooth stage, or replace harmonic tropicalizations by analytic moments. The other GR source, arXiv:2503.05513, concerns strict real-PL semipositivity and its curve comparison; it is not used to infer weak smoothness of arbitrary lisse functions. Berkovich 1993, Proposition 1.5.5, is the boundary-composition premise as restated in GR 2.11.

\section{Auxiliary arguments, including smoothness and positivity}

\subsection{Countable extraction with extended-real limits}

Let \(Z\) be second countable and let \((a_i)_{i\in I}\) be a directed decreasing net of finite continuous real functions with pointwise infimum \(a:Z\to\mathbb R\cup\{-\infty\}\). Then there are indices \(i_j\le i_{j+1}\) with \(a_{i_j}\downarrow a\) at every point.

For each rational \(q\), the open sets \(\{a_i<q\}\) cover the open set \(\{a<q\}\). Every open subspace of a second-countable space is Lindelof: given a cover, select for each member of a countable base contained in one cover member one such member. Thus countably many indices suffice for each rational \(q\). Enumerate all selected indices and take successively directed upper bounds. If \(a(z)\) is finite, every rational \(q>a(z)\) eventually bounds the selected values at \(z\); if \(a(z)=-\infty\), every rational does. The reverse inequality follows from \(a_i\ge a\). This is extraction by values, not a claim that the original index set has a countable cofinal subset.

If \(K\) has a countable dense subset, each strict affinoid algebra is separable in its Banach norm. Indeed, polynomials in finitely many variables with coefficients from a countable dense subfield are dense in a Tate algebra, and their images are dense in a quotient. For a norm-dense sequence \((b_n)\), evaluations \(x\mapsto |b_n(x)|\) generate the Berkovich topology: the inequality
\[
\big||b(x)|-|c(x)|\big|\le |(b-c)(x)|\le\|b-c\|
\]
gives uniform approximation of every evaluation. The countable evaluations separate points and embed the compact spectrum into a countable product of bounded real intervals. Hence the spectrum, and every ordinary open subset of it, is second countable. This proof is local and does not require global countability of an analytic space.

\subsection{Compact domination without continuity of the target}

If \(a_i\downarrow a\), the \(a_i\) are continuous, \(A\) is compact, and a finite continuous \(b\) satisfies \(b>a\) on \(A\), then some index satisfies \(a_i<b\) on all of \(A\). The open sets \(\{a_i<b\}\) cover \(A\). A finite subcover and one directed upper bound give the conclusion. Additional finitely many index or point-value requirements can be imposed by taking another upper bound. This argument works when \(a\) is discontinuous or minus infinity and is not Dini's theorem.

\subsection{A regularized maximum and classical positivity}

Choose an even nonnegative smooth probability density \(\rho\) supported in \((-1,1)\). For \(\delta>0\), set
\[
\theta_\delta(t)=\int_{-1}^1 |t+\delta s|\rho(s)\,ds,
\qquad M_\delta(a,b)=\frac{a+b+\theta_\delta(a-b)}2.
\]
Convolution makes \(\theta_\delta\) smooth and convex. Its derivative lies in \([-1,1]\). Jensen's inequality and the triangle inequality imply
\[
\max(a,b)\le M_\delta(a,b)\le\max(a,b)+\delta/2.
\]
It equals the maximum when \(|a-b|\ge\delta\), and is convex and nondecreasing in each variable. These facts prove all bounds used below.

If \(g_1,\ldots,g_r\) are finite continuous classical psh functions and \(F\) is a finite smooth convex function on an open convex region containing their values, nondecreasing in each variable, then \(F(g_1,\ldots,g_r)\) is classically psh. Supporting affine functions to \(F\) have nonnegative coefficients. Their finite maxima evaluated on the \(g_i\) are classically psh by GR 4.2. The directed net of finite maxima converges pointwise to the finite continuous composition. GR Theorem 6.5 applies. The argument also applies to a convex nondecreasing function of \(\log|P|\) which is constant on a negative half-line: its constant supporting line controls zeros, so the composition remains finite and continuous there. This proves classical positivity of every regularized maximum used below; it does not follow from its smoothness label.

\subsection{What the curve approximants actually look like}

We use GJR 15.11 as the published regularization premise. Its proof has a more precise local description which is needed for pullbacks and current positivity. After a finite Galois extension, an individual lisse target is represented on a finite skeleton. At a vertex with outgoing slopes \(s_e\) and weights \(w_e>0\), put
\[
\mu=\sum_e w_es_e>0,\qquad
g_e=s_e-\frac{\mu}{\sum_e w_e}.
\]
The star function with value \(F(x)\) and outgoing slopes \(g_e\) is harmonic, lies strictly below \(F\) off \(x\), and agrees with it at \(x\). A sufficiently small \(M_\delta(G+\epsilon,F)\), with \(\delta<\epsilon\), equals \(G+\epsilon\) near the vertex and equals \(F\) near the outer ends. These are equalities on ordinary neighborhoods, so the patched functions agree on the ordinary overlaps. Arbitrary real subdivisions in this argument are not asserted to come from strict formal models.

Off these small stars the function is harmonic. Near a vertex it is harmonic plus a constant. On each remaining annulus it is a smooth convex function \(\Phi(t)\) of an invertible logarithmic annulus coordinate \(t\). Harmonic regions are weakly smooth because GJR 15.7 and 7.5 write their real harmonic potential as a smooth affine function of finitely many integer harmonic coordinates. Annular regions are weakly smooth using the smooth tropicalization \(t\). The finite Galois construction can be made invariant by averaging its finitely many translates: averages retain weak smoothness, classical positivity, approximation bounds and decreasing order. Proposition 14.2 then descends the weakly smooth functions. The harmonic pieces are subharmonic, annular convex pieces are subharmonic, and ordinary sheafiness gives classical positivity, using GR 3.4.1 and 4.1.2. The identically minus-infinity component is handled separately by constants \(-j\).

For clarity, the local curvature check is explicit. Harmonic pieces have \(d'd''g=0\). Annular pieces have
\[
d'd''g=\Phi''(t)\,d't\wedge d''t,\qquad\Phi''\ge0.
\]
Convex compositions and regularized maxima have positive semidefinite Hessians in these harmonic coordinates. Under an analytic morphism, each integer harmonic coordinate pulls back to an integer harmonic coordinate by GJR 7.6, so the same formulas hold with \(t\) or the harmonic tuple composed with the morphism. A real harmonic coordinate is represented by an affine function of an integer harmonic tuple; its Hessian is zero. Thus these pullbacks are weakly smooth with positive curvature, rather than merely being assigned a smoothness name.

Here is the current interpretation whenever the target is Hausdorff, boundaryless, good strict, and pure of dimension \(d\). A positive semidefinite Hessian gives a strongly positive \((1,1)\)-form: diagonalize the matrix to obtain a sum of nonnegative multiples of \(\ell\wedge J\ell\). Wedge with a positive compactly supported test form of bidegree \((d-1,d-1)\); the Lagerberg positivity cones in that degree coincide by CD 1.2.10, so it is enough to use strongly positive tests. On a common harmonic tropical chart the resulting top-degree coefficients are nonnegative. Integration uses the nonnegative tropical multiplicities of GJR and is nonnegative. Localization over the compact support proves
\[
\langle d'd''[g],\eta\rangle=\int d'd''g\wedge\eta\ge0.
\]
Stokes in GJR 11--12 justifies the equality. If the formula was obtained only after a finite Galois extension, pull back \(\eta\) and use GJR 14.5(1): the integrals are equal, with no degree factor. Finite averaging is harmless. This proves positive strong curvature currents and hence positive CLD currents by restricting the test space. It is a proof of positivity for the constructed functions, not a general equivalence between classical and current positivity for all weakly smooth functions.

These facts remain true after every analytic extension \(L/K\). The coordinates pull back by GJR 8.2(4). A finite Galois auxiliary extension can be composited with \(L\), giving a finite Galois extension of \(L\); the same integral argument applies. Classical positivity also follows independently from GR 4.2(5). No extraction of a new sequence over \(L\) is needed for the pullback of a sequence already constructed over \(K\).

\section{A sequence theorem for singular curves, including several branches}

\textbf{Theorem A.} Let \(K\) be complete, perfect, nontrivially valued, with a countable dense subset. Let \(C\) be a reduced, Hausdorff, good, strictly \(K\)-analytic curve with empty analytic boundary. For every ordinary open \(D\subset C\), every classical psh \(u:D\to\mathbb R\cup\{-\infty\}\), and every \(x\in D\), there is one ordinary open neighborhood \(V\subset D\) of \(x\) and a sequence of finite weakly smooth functions \(v_j\) on \(V\) such that \(v_j\) is classically psh, \(d'd''[v_j]\) is a positive strong current, and \(v_j\downarrow u|_V\) pointwise. The same sequence pulls back with these properties after every analytic extension. If \(u\) is finite continuous, the convergence is locally uniform. This theorem permits singularities with several branches, as well as unibranch singularities. No residue-field algebraicity or Mumford hypothesis is used.

\textbf{Proof.} On the nonsingular locus, normalization is an isomorphism and is rig-smooth over the perfect field. Shrink to a relatively compact ordinary neighborhood in a smooth normalization component, apply GJR 15.11, handle a minus-infinity component by constants, and use 2.1. Its positivity and pullbacks are covered by 2.4.

Consider a singular point \(x\), necessarily rigid. By GR Lemma 3.9, shrink a connected neighborhood so that it has no other singularity and the finite normalization
\[
\nu:\coprod_{i=1}^r U_i\longrightarrow U
\]
has one preimage \(x_i\) of \(x\) on each connected smooth branch \(U_i\), and is an isomorphism away from these points. Choose the \(U_i\) relatively compact in slightly larger smooth neighborhoods on which \(u\circ\nu\) is defined, and inside a normalized strict affinoid neighborhood. These are second countable by 2.1. Such a choice is possible by finiteness, ordinary local compactness, and the local description in GR 3.9. Choose a compact strict affinoid neighborhood \(Y\subset U\) of \(x\), whose ordinary interior \(V\) contains \(x\). Let
\[
A_i=\nu^{-1}(Y)\cap U_i,\qquad A=\coprod_i A_i.
\]
These are compact. Write \(u_i=u\circ\nu|_{U_i}\) and \(a=u(x)=u_i(x_i)\). On each branch, GJR and 2.1 supply a decreasing sequence \(s_{i,j}\) of finite weakly smooth classical psh functions, with the curvature descriptions of 2.4, converging pointwise to \(u_i\). If \(u_i\equiv-\infty\), use \(s_{i,j}=-j\).

Every weakly smooth function is locally constant at a rigid point of a curve. Indeed, its local preform uses finitely many integer harmonic coordinates. Each is strict PL. At a rigid point a strict affinoid G-cover has a member containing an ordinary neighborhood of that point, as recorded in the proof of GR 3.9(6). Invertible analytic absolute values are locally constant there. Each coordinate, and therefore the smooth function of them, is locally constant. In the singular descent below we only use this assertion on the smooth normalization branches.

Take \(\epsilon_0=1\) and \(\epsilon_j=8^{-j}\). Upper semicontinuity bounds \(u_i\) above on the compact closures of the \(U_i\); choose one constant \(h_{i,0}=H\) for all branches with \(H\ge u_i+1\). We construct functions \(h_{i,j}\) on the fixed \(U_i\), satisfying \(h_{i,j}\ge u_i+\epsilon_j\), which decrease on \(A_i\) and have exactly the same value at all \(x_i\).

Suppose first that \(a\) is finite. Choose a sufficiently late term \(g_i\) of the branch sequence so that
\[
g_i\le s_{i,j},\qquad g_i<h_{i,j-1}-4\epsilon_j\ \text{on }A_i,
\qquad a\le b_i:=g_i(x_i)<a+\epsilon_j.
\]
The compact inequality is available by 2.2 because \(h_{i,j-1}\ge u_i+\epsilon_{j-1}\) and \(4\epsilon_j<\epsilon_{j-1}\). The point-value inequality follows from convergence. Put
\[
b=\max_i b_i,\qquad h_{i,j}=g_i+(b-b_i)+\epsilon_j.
\]
Then \(0\le b-b_i<\epsilon_j\), all \(h_{i,j}(x_i)=b+\epsilon_j\), and
\[
u_i+\epsilon_j\le h_{i,j}\le s_{i,j}+2\epsilon_j,
\qquad h_{i,j}<h_{i,j-1}-2\epsilon_j\ \text{on }A_i.
\]
Addition of constants preserves the smoothness and both positivity assertions. The common value and local constancy give equality on full neighborhoods of the normalization preimages, not just at isolated points.

Suppose next that \(a=-\infty\). Choose a finite number
\[
c_j\le-j,\qquad c_j<\min_A h_{j-1}-4\epsilon_j,
\]
where \(h_{j-1}\) denotes the continuous function on the disjoint union of branches. Choose a late \(g_i\) with
\[
g_i\le s_{i,j},\quad g_i<h_{i,j-1}-4\epsilon_j\ \text{on }A_i,
\quad g_i(x_i)<c_j-\epsilon_j.
\]
Define
\[
h_{i,j}=M_{\epsilon_j}(g_i,c_j)+\epsilon_j.
\]
It equals \(c_j+\epsilon_j\) near every \(x_i\). Section 2.3 proves classical positivity; the convex harmonic-chart formulas in 2.4 prove weak smoothness and current positivity. The bounds are
\[
u_i+\epsilon_j\le h_{i,j}\le\max(s_{i,j},c_j)+3\epsilon_j/2,
\qquad h_{i,j}<h_{i,j-1}-5\epsilon_j/2\ \text{on }A_i.
\]
They also cover branches identically minus infinity. Since \(c_j\to-\infty\), they prove pointwise convergence to \(u_i\).

In either case, the branch functions descend to finite continuous functions \(v_j\) on \(V\). To check ordinary smoothness-class descent, away from \(x\) use the normalization isomorphism. Near \(x\), choose the open neighborhoods on which every \(h_{i,j}\) has the common constant value. The finite normalization is closed and a quotient; the complement of the image of the complementary closed set gives an ordinary neighborhood of \(x\) on which \(v_j\) is constant. Thus \(v_j\) is weakly smooth in an actual neighborhood of the singularity and all chart overlaps agree. Finite surjective classical descent, GR 4.2(6), proves that \(v_j\) is classically psh. Alternatively, it is constant at the singularity and classically psh on the normalization-isomorphic complement. The curvature is zero on that constant neighborhood and has the positive formulas of 2.4 elsewhere, proving positive strong currents on \(V\).

The inequalities on the fixed compact \(A\) imply decrease on the fixed ordinary \(V\); the pointwise bounds imply its limit is \(u\), including its minus-infinity values. Section 2.4 proves scalar compatibility. On every compact subset where the target is finite continuous, ordinary Dini for this sequence gives uniform convergence. No uniform convergence is claimed for arbitrary singular targets. This proves the theorem.

The cusp \(y^2=x^3\) over a characteristic-zero field as in the theorem is a concrete unibranch example. A reduced nodal curve illustrates why independently choosing normalization approximants is insufficient: their values over the node must be aligned. The finite-branch construction above performs that alignment on one fixed neighborhood. It does not assert general weak-smooth descent along an arbitrary finite morphism.

\section{Noninvariant higher-dimensional sector on saturated proper-fibre domains}

\textbf{Theorem B.} Let \(K\) be as in Theorem A, and let \(C\) be a curve as there. Let \(f:X\to C\) be a proper surjective morphism with geometrically connected nonempty analytic fibres, where \(X\) is Hausdorff, good, strictly analytic, and pure of dimension \(d\ge2\). Suppose \(f\) has analytic sections locally on an ordinary open cover of \(C\). For every ordinary open \(D\subset C\) and every arbitrary classical psh function \(u\) on the \textbf{saturated domain} \(f^{-1}(D)\), every \(x\in f^{-1}(D)\) has a fixed ordinary open neighborhood \(f^{-1}(V)\), with \(f(x)\in V\subset D\), and finite weakly smooth functions
\[
w_j=v_j\circ f\quad\text{on }f^{-1}(V)
\]
which are classically psh, have positive strong curvature currents, and decrease pointwise to \(u\). For finite continuous \(u\), convergence is locally uniform. The same sequence works by pullback after every analytic extension. Fibres need not be smooth or normal.

\textbf{Proof of fibre constancy.} For \(c\in D\), pass to \(\mathcal H(c)\). GR scalar stability and pullback show that the restriction of \(u\) to the proper connected fibre \(X_c\) is classically psh. The fibre is compact and has empty analytic boundary. If it is not identically minus infinity, upper semicontinuity gives a finite attained maximum \(M\). The maximum set \(\{u=M\}\) is nonempty and closed. GR Theorem 6.1 makes it open, because every fibre point is an analytic interior point. Connectedness makes it the whole fibre. If the function is identically minus infinity, constancy already holds. Thus \(u\) has one value on each fibre, including fibres over nonrigid base points. This argument does not use density of rational points.

There is a unique function \(a:D\to\mathbb R\cup\{-\infty\}\) with \(u=a\circ f\). On an ordinary open admitting a section \(s\), \(a=u\circ s\), so it is upper semicontinuous and classically psh by GR pullback. On overlaps these functions are equal by fibre constancy, so ordinary sheafiness gives a classical psh function on \(D\). Apply Theorem A to \(a\) at \(f(x)\) and pull back its sequence.

The explicit local preforms are \((h\circ f,\alpha)\) for a harmonic chart \((h,\alpha)\) for \(v_j\). GJR 7.6 makes the coordinates harmonic. The common-chart equality on overlaps follows by pulling back the joint chart \((h,h')\), as in Definitions 10.3 and Proposition 10.9. The positive Hessian/zero-curvature representations, including the constant neighborhoods over singular base points, pull back exactly as in 2.4. Hence the strong currents are positive. Since \(f\) is proper and \(C\) has empty boundary, GR 2.11 implies \(X\) and \(f^{-1}(V)\) have empty analytic boundary, so these currents are defined without boundary residues.

For any analytic \(L/K\), let \(p:X_L\to X\). On the one fixed ordinary domain \(p^{-1}f^{-1}(V)=f_L^{-1}(V_L)\), the sequence is \(w_j\circ p\). It remains decreasing, and its value at a point \(z\) tends to \(u(p(z))\). Harmonic pullback, the scalar positivity argument in 2.4, and GR 4.2(5) prove all other assertions. This constructs one base-field sequence and then extends it; it does not attempt descent of arbitrarily chosen extension-field approximants. If \(u\) is finite continuous, a local section makes \(a\) finite continuous; compact images under \(f\) and the local uniform convergence of \(v_j\) give local uniform convergence of \(w_j\).

Examples include \(X=F^{\mathrm{an}}\times C\) for a geometrically connected proper algebraic \(F\) with a \(K\)-point, provided the displayed space is pure-dimensional. The point gives a section and projectivity supplies actual examples such as \(F=\mathbb P^m\), or a geometrically connected reduced nodal projective curve. The latter gives a singular surface when \(C\) is a curve. For \(C=\mathbb A^{1,\mathrm{an}}\), the function \(\log|T-a|\) pulled back from the base, with \(a\ne0\), is an arbitrary-class example with a genuine pole; on a torus open its base potential is not generally maximal-compact-torus invariant. The theorem assumes no torus invariance of any starting function.

\textbf{The domain quantifier matters.} Theorem B regularizes every classical psh function on each saturated domain \(f^{-1}(D)\). It does not regularize every sheaf germ on every ordinary open subset of \(X\). For instance, on an affine chart of \(\mathbb P^m\times C\), \(\log|Z_1|\) is classical psh and varies along the affine fibres; the compact-fibre argument does not apply there. Properness of \(X\) or of its fibres therefore cannot be advertised as a characterization of the original local sheaf property.

\section{A sharp analytic-boundary family}

\textbf{Theorem C.} For a complete nontrivially valued \(K\), let
\[
E=\mathcal M(K\langle T\rangle),\qquad \xi=\text{its radius-one Gauss point},\qquad
b(\xi)=1,\quad b(z)=0\ (z\ne\xi).
\]
This is an upper semicontinuous classical psh function on the good strict rig-smooth closed disc. If \(\widetilde K\) is uncountable, on no ordinary neighborhood of \(\xi\) is \(b\) a pointwise limit of a decreasing sequence of finite continuous functions; indeed, the argument below excludes any pointwise convergent sequence of continuous functions. If \(\widetilde K\) is countable, there is an explicit decreasing sequence of finite weakly smooth classical psh functions on the whole fixed closed disc converging to \(b\), with positive curvature on its analytic interior. Thus residue countability is an exact threshold for this particular boundary-spike function, not a characterization for all psh functions on discs.

\textbf{Classical membership and the actual boundary.} The analytic boundary of \(E\) is \(\{\xi\}\). The function is upper semicontinuous since its strict sublevel set for \(0<a\le1\) is \(E\setminus\{\xi\}\). It is constant zero on \(\operatorname{Int}(E)\). GR Proposition 3.5(3) makes it subharmonic on \(E\), and GR 4.1.2 makes it classically psh. This supplies, in particular, the required tests after every analytic extension. It is not an interior spike on a boundaryless space.

\textbf{Uncountable-residue obstruction.} Every ordinary neighborhood \(N\) of \(\xi\) contains the rational points in all but finitely many residue classes of the unit disc. To prove this, choose a basic neighborhood specified by finitely many polynomial absolute-value conditions. Such conditions generate the spectrum topology: analytic functions are uniformly approximated by polynomials, and their norms at \(\xi\) are nonzero unless the function is zero. Normalize each nonzero polynomial by a coefficient of maximum absolute value. Its reduction is a nonzero polynomial in \(\widetilde K[T]\), which has only finitely many roots in \(\widetilde K\). For a lift \(a\in K^\circ\) of a residue class avoiding all these roots, each polynomial has precisely its Gauss absolute value at \(a\); hence \(a\in N\).

Suppose continuous functions \(q_j\) on a neighborhood \(N\) converge pointwise to \(b\). Since \(q_j(\xi)\to1\), discard finitely many terms to have \(q_j(\xi)>3/4\). Each \(N_j=\{q_j>3/4\}\cap N\) is an ordinary neighborhood of \(\xi\). Each excludes rational points from only finitely many residue classes by the preceding argument. Choose a residue class outside the countable union of exceptions, and one lift \(a\). Then \(q_j(a)>3/4\) for all remaining \(j\), contrary to \(b(a)=0\). Weakly smooth functions are continuous, so this is a fortiori a failure of the requested sequential regularization. It does not exclude decreasing nets.

\textbf{Countable-residue construction.} Enumerate the monic irreducible polynomials of \(\widetilde K[T]\) as \(p_1,p_2,\ldots\), and choose coefficientwise lifts \(P_i\in K^\circ[T]\) which are monic and have Gauss norm one. Choose a nonnegative smooth probability density \(\sigma\) supported in \((0,1)\), symmetric about \(1/2\). Set
\[
\eta(t)=2\int_t^1\sigma(s)\,ds,\qquad
\psi(t)=\int_0^t\eta(s)\,ds,
\]
where \(\sigma\) is zero outside \((0,1)\). Then \(\eta\) is smooth, nonnegative and nonincreasing; it is \(2\) for \(t\le0\) and zero for \(t\ge1\). The function \(\psi\) is smooth, increasing, concave, \(\psi(0)=0\), and \(\psi(t)=1\) for \(t\ge1\), because \(2\int_0^1 s\sigma(s)\,ds=1\). Define
\[
v_j(z)=1-\sum_{i=1}^j\psi\bigl(j[-\log|P_i(z)|]\bigr),
\]
with the summand defined as \(1\) at a zero of \(P_i\). This convention agrees on a neighborhood of every zero since \(\psi\) is identically one for arguments at least one.

At any point of \(E\) at most one \(|P_i|\) is less than one. Two distinct reductions are coprime, so lifting a Bezout identity would otherwise express \(1\) as terms of absolute value strictly less than one, plus a polynomial of Gauss norm strictly less than one. This is impossible. Consequently \(0\le v_j\le1\). Each existing summand increases with \(j\), and the new summand is nonnegative, so \(v_{j+1}\le v_j\).

If all \(|P_i(z)|=1\), every nonzero polynomial has its Gauss norm at \(z\): normalize it and factor its reduction into the \(p_i\); the difference between that polynomial and a corresponding product of lifts has Gauss norm less than one. Thus \(z=\xi\), by equality of polynomial seminorms. For every other point there is one \(i\) with \(|P_i(z)|<1\); its summand is eventually exactly one. Hence \(v_j(z)\to0\) there, while \(v_j(\xi)=1\) for all \(j\).

To check smoothness and positivity, put \(H_j(s)=1-\psi(-js)\). It is smooth, convex, and nondecreasing on \(\mathbb R\), and equals zero for \(s\le-1/j\). Each \(H_j(\log|P_i|)\) is locally constant near the zero set and elsewhere is a smooth function of an invertible analytic logarithm. It is therefore ordinary CLD smooth, hence weakly smooth through that smooth harmonic tropicalization. Classical positivity follows from the supporting-line argument in 2.3, including zeros. Since
\[
v_j=\sum_{i=1}^jH_j(\log|P_i|)-(j-1),
\]
the cone and constant properties give classical positivity. Its local Hessian is a sum of terms
\[
H_j''(\log|P_i|)\,d'\log|P_i|\wedge d''\log|P_i|\ge0.
\]
On the analytic interior the associated strong and CLD currents are positive by 2.4. We make no boundary-free integration-by-parts claim on the entire closed disc. All expressions pull back under every analytic extension, remain weakly smooth and classically psh, remain decreasing, and converge to the pullback \(b\circ p\). The limit on an extension is that pullback, not a newly defined spike at the extension-field Gauss point alone.

\textbf{Higher-dimensional failure.} For any \(m\ge1\), put \(X=E\times\mathbb P^{m,\mathrm{an}}\) and \(u=b\circ\operatorname{pr}_E\). This is an arbitrary-class psh function on a separated good strict space of pure dimension \(m+1\), by GR pullback. Fix a \(K\)-rational projective point \(p\) and \(x=(\xi,p)\). The section \(s(z)=(z,p)\) pulls any alleged sequence on an ordinary neighborhood of \(x\) back to a sequence on an ordinary neighborhood of \(\xi\), which was proved impossible when the residue field is uncountable. Projection is proper and has no relative boundary. Boundary composition in GR 2.11, with the disc locally separated, gives
\[
\partial X=\operatorname{pr}_E^{-1}(\partial E).
\]
Thus \(x\) is an actual analytic boundary point. A concrete field is \(K=\mathbb C((t))\) with its complete nontrivial \(t\)-adic valuation. The obstruction survives passage to any analytic extension with uncountable residue field; the original failure needs no scalar descent argument.

\section{Remaining original obligations and limits}

Theorems A--C do not characterize all higher-dimensional good strict spaces with the P02.3 sheaf property. The positive higher-dimensional theorem is restricted to functions defined on saturated proper-fibre domains over the indicated singular curves. It does not settle nonsaturated affine germs, even on those products. The full boundaryless higher-dimensional arbitrary-class problem, a sharp general singularity range, nonperfect fields, arbitrary nonseparable fields for general singular targets, nonseparated spaces, and general analytic boundary configurations remain unresolved here. The exact residue-countability threshold proved in Section 5 concerns one concrete boundary family. It supplies a necessary exclusion for a universal sequential statement over all fields and spaces, not a sufficient geometric characterization. None of our assertions identifies weakly smooth with ordinary CLD smooth, or derives current positivity just from class membership. There is no assertion that a general regularization net has a countable cofinal subnet.

\section{Reproducibility and evidence classification}

The input is retained in input.json and byte-identical input-original.json; assignment-original.txt and task-original.json retain its selected fields. The complete proof is this file; P02.3-paper.tex embeds the same body as a standalone paper. source-extracts/source-provenance.json records SHA-256 checks, literal source-page excerpts including complete selected published proofs, exact premises used, and their application conditions. PDF pages 80--81 of GJR were rendered and visually inspected to resolve the decisive net/sequence and subdivision wording. The TeX diagnostics are in build-1.stdout.txt, build-2.stdout.txt, build-1.log, build-2.log, and build-diagnostics.json. The existing TeX installation is used without installation.

Rerun the retained source audit with \texttt{python3 collect\_sources.py}, generate the standalone TeX with \texttt{python3 make\_paper.py}, and build twice with \texttt{python3 build\_paper.py}. \texttt{python3 audit\_artifacts.py} checks byte preservation, required status keys, proof-body synchronization, and recorded build diagnostics. These are artifact checks, not proof checking. No computational experiment or exhaustive finite computation is presented as evidence for the mathematics. The mathematical claims are draft proofs relying on the explicitly listed published foundations. No heavy computation, external service, private campaign files, other model, or subagent was used.
\end{document}
